\section{Pure trace on the constant-payoff domain}
\label{pt:sec:environment}\label{lf:app:pure-trace}

\subsection{Binary and finite branching}

We first record the finite-branch form of Axiom~\textup{(A8)} used throughout the proofs.  The result also makes
explicit where the other structural axioms enter the equivalence.

\begin{lemma}[Finite-branch extension]\label{pt:lem:finite-branch-extension}
Suppose Axioms~\textup{(A1)} and \textup{(A5)--(A7)} hold.  Then Axiom~\textup{(A8)} is equivalent to the
following property: for every $k\ge1$, every first-stage record channel
$P:A\to\D(\{1,\ldots,k\})$, every $q\in\D(A)$, and every
two ordered lists of continuation channels $(K_1,\ldots,K_k)$ and $(L_1,\ldots,L_k)$, with
$K_i,L_i:A\to\D(O\times R)$, if
\[
 (q_i,K_i)\succeq(q_i,L_i)
 \quad\text{for every reached }i,
\]
then
\begin{equation}
 \bigl(q,P\triangleright(K_1,\ldots,K_k)\bigr)
 \succeq
 \bigl(q,P\triangleright(L_1,\ldots,L_k)\bigr),
 \label{pt:eq:finite-branch-extension}
\end{equation}
strictly if one reached branch comparison is strict.
\end{lemma}

\begin{proof}
Assume first Axiom~\textup{(A8)}.  Put
\[
 m(i):=\sum_aq(a)P(i\mid a),
 \qquad
 I^+:=\{i\in\{1,\ldots,k\}:m(i)>0\}.
\]
If $m(i)=0$, then $P(i\mid a)=0$ for every $a\in\supp q$.  Consequently, changing the continuation after such a
record changes the compound channel only on action rows outside $\supp q$.  Apply Axiom~\textup{(A7)} in both
directions using the identity action processor: its Bayesian pushforwards agree on supported rows and allow
arbitrary completions elsewhere.  Continuations following unreached records are therefore immaterial.  We may
hold them fixed and enumerate the reached records as $I^+=\{i_1,\ldots,i_n\}$.

If $n=1$, the first-stage record is constant on $\supp q$.  Deleting or restoring that constant record is an
invertible record relabelling, up to rows outside $\supp q$; Axioms~\textup{(A6)} and \textup{(A7)} therefore
identify the compound comparison with the sole branch comparison.  Suppose henceforth that $n\ge2$.

For $j=0,\ldots,n$, define hybrid continuation channels by
\[
 H_i^j:=
 \begin{cases}
 K_i,&i\in\{i_1,\ldots,i_j\},\\
 L_i,&i\in I^+\setminus\{i_1,\ldots,i_j\},
 \end{cases}
\]
and use a common arbitrary continuation on $\{1,\ldots,k\}\setminus I^+$.  Write
$E_j:=P\triangleright(H_1^j,\ldots,H_k^j)$.  Fix $j\ge1$ and collapse the first-stage record to the binary channel
\[
 B_j(1\mid a):=P(i_j\mid a),
 \qquad
 B_j(2\mid a):=\sum_{i\ne i_j}P(i\mid a).
\]
Both binary records are reached.  Put all binary continuations on the common record set
$\{1,\ldots,k\}\times R$.  Pad the
two continuations that may differ by
\[
 \widetilde K_j(o,(i,r)\mid a):=\1_{\{i=i_j\}}K_{i_j}(o,r\mid a),
 \qquad
 \widetilde L_j(o,(i,r)\mid a):=\1_{\{i=i_j\}}L_{i_j}(o,r\mid a).
\]
On record $2$, collect all unchanged branches in the continuation channel
\begin{equation}
 M_j(o,(i,r)\mid a)
 :=\frac{P(i\mid a)}{B_j(2\mid a)}H_i^j(o,r\mid a),
 \qquad i\ne i_j,
 \label{pt:eq:complement-continuation}
\end{equation}
with zero probability on records $(i_j,r)$, whenever $B_j(2\mid a)>0$, and complete the other rows arbitrarily.
Since $H_i^j=H_i^{j-1}$ for $i\ne i_j$, this continuation is common to the two comparisons.
Binary record $1$ has probability
$m(i_j)$ and carries the conditional action lottery $q_{i_j}$; record $2$ carries the complementary conditional
lottery, and their mixture is $q$.  Appending or deleting the constant label $i_j$ is an invertible record
relabeling, so Axiom~\textup{(A6)} identifies the comparison of $\widetilde K_j$ and $\widetilde L_j$ with that of
$K_{i_j}$ and $L_{i_j}$.  Hence Axiom~\textup{(A8)} gives
\[
 \bigl(q,B_j\triangleright(\widetilde K_j,M_j)\bigr)
 \succeq
 \bigl(q,B_j\triangleright(\widetilde L_j,M_j)\bigr),
\]
strictly whenever the comparison at $q_{i_j}$ is strict.

The nested records $(1,(i_j,r))$ and $(2,(i,r))$ correspond bijectively to the flat records $(i,r)$ through
\[
 (1,(i_j,r))\longmapsto(i_j,r),
 \qquad
 (2,(i,r))\longmapsto(i,r)\quad(i\ne i_j).
\]
This relabelling is a payoff-preserving record processing in both directions.  Axiom~\textup{(A6)} therefore
makes the two binary compounds indifferent to $E_j$ and $E_{j-1}$, respectively.  Thus
$(q,E_j)\succeq(q,E_{j-1})$, strictly at any strict replacement.  Place the finite collection
$E_0,\ldots,E_n$ in one block environment.  Axiom~\textup{(A5)} transports each pairwise comparison to that
environment, and transitivity in Axiom~\textup{(A1)} yields $(q,E_n)\succeq(q,E_0)$, strictly if at least one
step is strict.  Restoring the immaterial null branches proves \eqref{pt:eq:finite-branch-extension}.

Conversely, assume the displayed finite-branch property, including its strict clause.  In the binary environment
of Axiom~\textup{(A8)}, reflexivity and duplication coherence make the common continuation weakly comparable to
itself, so the finite-branch property gives the forward implication in
\eqref{eq:recordwise-sure-thing}.  If the aggregate comparison in that equation held while
$(q_1,K)\not\succeq(q_1,L)$, completeness would give $(q_1,L)\succ(q_1,K)$.  Applying the strict
finite-branch property to the reversed lists would give the strict reverse aggregate comparison, a contradiction.  This
proves the converse implication and hence Axiom~\textup{(A8)}.
\end{proof}

\begin{corollary}[One-branch insertion]\label{pt:cor:one-branch-a8}
Under the hypotheses of Lemma~\ref{pt:lem:finite-branch-extension}, fix $q\in\D(A)$, a first-stage record channel
$P:A\to\D(\{1,\ldots,k\})$, a reached record $i^*$, and continuation channels
$K_1,\ldots,K_k,L_1,\ldots,L_k:A\to\D(O\times R)$ such that $K_i=L_i$ for every $i\ne i^*$.  Then
\[
 (q_{i^*},K_{i^*})\succeq(q_{i^*},L_{i^*})
 \quad\Longleftrightarrow\quad
 \bigl(q,P\triangleright(K_1,\ldots,K_k)\bigr)
 \succeq
 \bigl(q,P\triangleright(L_1,\ldots,L_k)\bigr),
\]
and the same equivalence holds for strict preference.
\end{corollary}

\begin{proof}
The forward implication follows from Lemma~\ref{pt:lem:finite-branch-extension}, using duplication coherence on
the common branches.  If the aggregate comparison held but the branch comparison failed, completeness would make
the reversed branch comparison strict; the strict part of the finite-branch extension would then give a strict
reverse aggregate comparison.  The strict equivalence follows by applying the weak equivalence in both directions.
\end{proof}

In the appendices, references to Axiom~\textup{(A8)} with more than two first-stage records mean its
finite-branch extension in Lemma~\ref{pt:lem:finite-branch-extension}; when only one branch changes, we cite
Corollary~\ref{pt:cor:one-branch-a8}.

The relevance axioms in the main text are stated inside fixed channels.  The
proofs below use their equivalent comparison-environment forms.

For $o\in O$, write $K_o:\{\ast\}\to\D(O\times\{\ast\})$ for the channel that
delivers $o$ for sure.  On a finite action set $A$, write
$K^{\mathrm{id}}_o:A\to\D(O\times A)$ for the channel that delivers $o$ and
reports the action, and $K^0_o:A\to\D(O\times\{\ast\})$ for the channel that
delivers $o$ and no record.

\begin{lemma}[Fixed-channel and environment relevance]\label{lem:relevance-bridge}
Given Axioms~\textup{(A1)}, \textup{(A5)}, \textup{(A6)},
\textup{(A7)}, and \textup{(A8)}, the following equivalences hold.
\begin{enumerate}[label=\textup{(\alph*)},leftmargin=*]
\item Axiom~\textup{(A3)} is equivalent to the existence of
  $o^+,o^-\in O$ such that
  \[
    (\delta_\ast,K_{o^+})\succ(\delta_\ast,K_{o^-}).
  \]
\item Axiom~\textup{(A4)} is equivalent to the following condition: there is
  an outcome $o_\ast\in O$ such that, for every finite $A$ with $|A|\ge2$
  and every full-support $q\in\D(A)$,
  \[
    (q,K^{\mathrm{id}}_{o_\ast})\succ(q,K^0_{o_\ast}).
  \]
\end{enumerate}
\end{lemma}

\begin{proof}
For part \textup{(a)}, apply action processing to either pure lottery in the
two-action channel of Axiom~\textup{(A3)}.  Collapsing its support to
a singleton gives $K_{o^+}$ or $K_{o^-}$.  Conversely, the singleton can be
processed back to the corresponding action; the unrestricted row at the
unreached action completes the pushforward to the original channel.  Thus
each pure lottery is indifferent to the corresponding singleton experiment.
Weak order and block coherence transport the strict comparison in either
direction.

For part \textup{(b)}, let $o_\ast$ witness Axiom~\textup{(A4)} and first
restrict its two lotteries to their supports.  Action processing in both
directions identifies the first with fair binary full revelation and the
second with a fair binary channel whose two rows are the same.  Record
processing in both directions identifies the latter channel with silence.
Hence the fixed-channel axiom is equivalent to
\begin{equation}
  ((\tfrac12,\tfrac12),K^{\mathrm{id}}_{o_\ast})
  \succ
  ((\tfrac12,\tfrac12),K^0_{o_\ast})
  \label{eq:bridge-fair-binary}
\end{equation}
on a binary action set.

Now let $p$ have full support on a binary set.  Choose
$0<\varepsilon\le 2\min_i p_i$ and an interim record $y$ with
$P(y\mid i)=\varepsilon/(2p_i)$.  The branch has probability $\varepsilon$
and posterior $(\tfrac12,\tfrac12)$.  Compare two continuation plans: after that branch the
first reveals the action and the second is silent; after every other branch
both are silent.  Equation~\eqref{eq:bridge-fair-binary} and
Corollary~\ref{pt:cor:one-branch-a8} give a strict comparison of the two compounds.  Full revelation
can be garbled into the first compound, and the second compound can be
garbled into silence.  Record processing and transitivity therefore give
$(p,K^{\mathrm{id}}_{o_\ast})\succ(p,K^0_{o_\ast})$.

Finally, let $q$ have full support on any finite $A$ with $|A|\ge2$, and
partition $A$ into two nonempty cells.  The channel that reports only the cell
is, by action processing in both directions, indifferent to binary full
revelation under the induced binary prior; silence on $A$ is similarly
indifferent to binary silence.  The cell-reporting channel is therefore
strictly preferred to silence, and full revelation is weakly preferred to
the cell report by record processing.  This proves the displayed
comparison for every $q$ and $A$ at $o_\ast$.  The converse follows by taking
a fair binary prior and reversing the support identifications used above.
\end{proof}

Whenever Axiom~\textup{(A4)} is imposed below, fix one of its witnessing
outcomes and denote it by $o_\ast$.  Every channel between nonempty finite sets
$P:A\to\D(X)$ has the constant-payoff lift
\begin{equation}
K_P(o,x\mid a):=\1_{\{o=o_\ast\}}P(x\mid a).
\label{pt:eq:constant-payoff-lift}
\end{equation}
Throughout this appendix, comparisons of pure record channels mean comparisons
of these lifts in the model already defined in the main text:
\[
q\succeq_Pq'\quad:\Longleftrightarrow\quad q\succeq_{K_P}q'.
\]
We similarly write $(q,P)\succeq(r,Q)$ for the main-text two-block
comparison of $(q,K_P)$ and $(r,K_Q)$.
Likewise, $P\sqcup Q$, block-supported lotteries, action and record
processing, ordered lists of continuation channels, and compounds
$P_1\triangleright(Q_1,\ldots,Q_k)$ are shorthand for the
corresponding main-text constructions after applying
\eqref{pt:eq:constant-payoff-lift}.  The lift commutes with each construction
up to canonical tagged or singleton-record bijections; those bijections are
neutral by Axiom~\textup{(A6)} applied in both directions.  No new primitive
relation or behavioural condition is introduced here.

Two pieces of notation will be used repeatedly.  If
$S:A\to\Delta(B)$ is an action processor, let $S^qP:B\to\Delta(X)$ denote any
Bayesian pushforward satisfying
\[
 (qS)(b)(S^qP)(x\mid b)
 =\sum_a q(a)S(b\mid a)P(x\mid a).
\]
Its rows at zero-mass reports are arbitrary, exactly as in
\eqref{eq:action-processing}.  For channels $P:A\to\Delta(X)$ and
$Q:B\to\Delta(Y)$, their independent product is
\[
 (P\otimes Q)((x,y)\mid(a,b)):=P(x\mid a)Q(y\mid b).
\]

For $q\in\D(A)$ and $P:A\to\D(X)$, write
\[
m_{qP}(x):=\sum_a q(a)P(x\mid a),
\qquad
r_x^{qP}(a):=\frac{q(a)P(x\mid a)}{m_{qP}(x)}
\quad\text{when }m_{qP}(x)>0,
\]
and define the finite posterior law
\[
\mu_{qP}:=
\sum_{x:m_{qP}(x)>0}m_{qP}(x)\,\delta_{r_x^{qP}}.
\]
Let $\Id_A$ be full revelation and $U_A$ the one-record uninformative
channel.  With logarithms in the fixed base used in the main text, put
\[
I_{qP}(A;X)
:=\Sh(m_{qP})-\sum_a q(a)\Sh(P(\cdot\mid a))
=\Sh(q)-\int_{\D(A)}\Sh(r)\,d\mu_{qP}(r).
\]

\begin{proposition}[Pure-trace lemma]\label{pt:prop:main}
Suppose the preference family satisfies main-text Axioms
\textup{(A1)}, \textup{(A2)}, and \textup{(A4)--(A8)}.  Then, for every record channel between
nonempty finite sets
$P:A\to\D(X)$ and every $q,q'\in\D(A)$,
\[
q\succeq_Pq'
\quad\Longleftrightarrow\quad
I_{qP}(A;X)\ge I_{q'P}(A;X).
\]
The same numerical scale represents every block-supported comparison: for
every nonempty finite family $\{P_i:A_i\to\D(X_i)\}_{i\in J}$ of channels
between nonempty finite sets, distinct $i,j\in J$,
and priors $q_i\in\D(A_i)$ and $q_j\in\D(A_j)$,
\[
q_i^i\succeq_{\bigsqcup_{k\in J}P_k}q_j^j
\quad\Longleftrightarrow\quad
I_{q_iP_i}(A_i;X_i)\ge I_{q_jP_j}(A_j;X_j).
\]
\end{proposition}

\subsection{Proof}

Assume the hypotheses of Proposition~\ref{pt:prop:main}.  Every appeal below
to weak order, continuity, block coherence, record processing, action
processing, the finite-branch extension of the recordwise sure-thing principle, or positive trace orientation is an
appeal to Axiom \textup{(A1)}, \textup{(A2)}, \textup{(A5)}, \textup{(A6)},
\textup{(A7)}, \textup{(A8)}, or \textup{(A4)}, respectively, on the
constant-payoff lifts.  Thus the proof has no auxiliary axiom system.

\paragraph{How to read the proof.}
The argument has five interfaces.  Each stage ends with exactly the object
needed by the next one:
\begin{enumerate}[label=\textup{Stage \arabic*:},leftmargin=*]
\item replace channels by Bayes-plausible posterior laws, derive public-mixture
  independence, and cardinalise each fixed-prior quotient using only the
  segment closedness supplied by primitive continuity;
\item make relabelling, undetectable distinctions, and restriction to the
  positive support exact consequences of action processing;
\item transport those affine values through independent products and into one
  cross-prior block scale, allowing provisionally a bilinear interaction;
\item identify reached-branch coefficients, prove their cocycle, and align
  the scales on nested support faces;
\item compare product and sequential decompositions, eliminate the
  interaction, and invoke Faddeev's recursion to obtain Shannon entropy and
  hence mutual information.
\end{enumerate}
The product-coefficient and branch tangent-space calculations are the two
technical subarguments.  A reader interested in the logical spine may use
their lemma statements as interfaces.  Product calibration is
placed before branch calibration because this makes the scalar algebra easier
to read.  The early full-revelation normalisation and exact relabelling lemma
below supply the coherence needed by that order of presentation.

\subsubsection{Stage 1: posterior-law quotient and affine fibre values}

For a finite action set $A$, let $\mathsf X_A$ be the set of finitely
supported probability measures on $\Delta(A)$, and let
\[
 b(\mu):=\int_{\Delta(A)}s\,d\mu(s),
 \qquad
 \mathcal M_q:=\{\mu\in\mathsf X_A:b(\mu)=q\}.
\]
Thus $\mathcal M_q$ is the mixture set of finite Bayes-plausible posterior
laws with prior $q$, and $\mu_{qP}\in\mathcal M_q$ for every finite channel
$P$.

\begin{lemma}[Coherent comparisons across blocks]\label{pt:lem:blockcoh}
The relation on prior--channel pairs defined above is complete and transitive.
It is unaffected by adding or deleting blocks other than the two being
compared.  In particular, for fixed $q$ the induced relation on channels is a
weak order, and
\[
 q\succeq_P r
 \quad\Longleftrightarrow\quad
 (q,P)\succeq(r,P).
\]
For fixed alphabets, its graph is closed under entrywise convergence of the
two channels and convergence of the two priors.
\end{lemma}

\begin{proof}
Completeness is Axiom~\textup{(A1)} in the two-block environment.  To prove
transitivity, suppose
$(q_1,P_1)\succeq(q_2,P_2)$ and
$(q_2,P_2)\succeq(q_3,P_3)$.  Put the three pairs in the common environment
$P_1\sqcup P_2\sqcup P_3$.  Axiom~\textup{(A5)} transports the two
hypotheses to that environment, Axiom~\textup{(A1)} composes them there, and
Axiom~\textup{(A5)} deletes the middle block.  Applying Axiom~\textup{(A5)}
with the ordered pair of block labels reversed gives
\[
 (q_2,P_2)\succeq(q_1,P_1)
 \quad\Longleftrightarrow\quad
 q_2^1\succeq_{P_1\sqcup P_2}q_1^0,
\]
so reversal is also canonical.  The displayed within-channel equivalence is
the duplication clause of Axiom~\textup{(A5)}.  Finally, a block channel is a
continuous affine function of its component channels.  Closedness on fixed
alphabets is therefore exactly Axiom~\textup{(A2)} applied to the common block
environment.
\end{proof}

\begin{lemma}[Posterior-law quotient]\label{pt:lem:blackwell}\label{pt:lem:plsuff}
Fix a full-support $q\in\Delta(A)$.  If finite channels
$P:A\to\Delta(X)$ and $Q:A\to\Delta(Y)$ satisfy
\[
 \mu_{qP}=\mu_{qQ},
\]
then they are mutually obtainable by record post-processing.  Consequently
\[
 (q,P)\sim(q,Q),
\]
and either channel may be substituted for the other in every block
comparison at prior $q$.  Hence the fixed-$q$ order of channels descends to a
well-defined weak order $\succeq_q$ on $\mathcal M_q$.
\end{lemma}

\begin{proof}
We give the finite matching argument because it also handles repeated and
zero-probability records.  Put
\[
 m_P(x):=\sum_aq(a)P(x\mid a),
 \qquad
 m_Q(y):=\sum_aq(a)Q(y\mid a).
\]
For every posterior $s$ occurring with positive probability, let
\[
 X_s:=\{x:m_P(x)>0,\ r_x^{qP}=s\},
 \qquad
 Y_s:=\{y:m_Q(y)>0,\ r_y^{qQ}=s\}.
\]
Equality of posterior laws gives the common positive mass
\[
 \lambda_s:=\sum_{x\in X_s}m_P(x)
            =\sum_{y\in Y_s}m_Q(y).
\]
For $x\in X_s$, define a stochastic record kernel by
\[
 T(y\mid x):=
 \begin{cases}
   m_Q(y)/\lambda_s,&y\in Y_s,\\
   0,&y\notin Y_s.
 \end{cases}
\]
Choose arbitrary stochastic rows at records $x$ with $m_P(x)=0$.  Full support
of $q$ implies that such a record has $P(x\mid a)=0$ for every $a$.

If $y\in Y_s$, Bayes' rule gives
\[
 \sum_{x\in X_s}P(x\mid a)
 =\frac{\lambda_s s(a)}{q(a)},
 \qquad
 Q(y\mid a)=\frac{m_Q(y)s(a)}{q(a)}.
\]
It follows that
\[
 (PT)(y\mid a)
 =\sum_{x\in X_s}P(x\mid a)\frac{m_Q(y)}{\lambda_s}
 =Q(y\mid a).
\]
Both sides vanish when $m_Q(y)=0$, again by full support.  Thus $PT=Q$.
Interchanging $P$ and $Q$ constructs $T'$ with $QT'=P$.

Axiom~\textup{(A6)} applied to $T$ and $T'$ yields the two weak comparisons;
Lemma~\ref{pt:lem:blockcoh} puts the reverse comparison into the same block
order and gives indifference.  If $P,P'$ have the same posterior law, place
$(q,P),(q,P')$ and the two objects being compared in one finite block
environment.  Transitivity then permits replacement of $P$ by $P'$ on either
side.  The quotient order is consequently independent of the implementing
channel and inherits completeness and transitivity from
Lemma~\ref{pt:lem:blockcoh}.
\end{proof}

Every element of $\mathcal M_q$ is implementable when $q$ has full support.
Indeed, if
\[
 \mu=\sum_{k=1}^n\lambda_k\delta_{s_k},
 \qquad \lambda_k>0,
 \qquad \sum_k\lambda_ks_k=q,
\]
then
\begin{equation}
 C_\mu(k\mid a):=\frac{\lambda_ks_k(a)}{q(a)}
 \label{pt:eq:canonical-implementation}
\end{equation}
defines a channel and $\mu_{qC_\mu}=\mu$.

For channels $P:A\to\Delta(X)$ and $R:A\to\Delta(Z)$ and
$t\in[0,1]$, let $tP\oplus(1-t)R$ be the publicly tagged mixture: it
reports which component was selected as well as that component's record.  Its
posterior law is
\begin{equation}
 \mu_{q,tP\oplus(1-t)R}
 =t\mu_{qP}+(1-t)\mu_{qR}.
 \label{pt:eq:public-posterior-mixture}
\end{equation}

\begin{lemma}[Public-mixture independence]\label{pt:lem:publiccoin}
Let $q$ have full support.  For all finite channels $P,Q,R$ on the action set
$A$ and every $t\in(0,1)$,
\[
 (q,P)\succeq(q,Q)
 \quad\Longleftrightarrow\quad
 (q,tP\oplus(1-t)R)
 \succeq
 (q,tQ\oplus(1-t)R).
\]
Equivalently, the quotient order on $\mathcal M_q$ satisfies mixture
independence.
\end{lemma}

\begin{proof}
If necessary, embed the record alphabets of $P,Q,R$ into one tagged disjoint
union and fill the unused coordinates with zeros.  This does not change any
posterior law, so Lemma~\ref{pt:lem:plsuff} makes the padding neutral.

Let $B:A\to\Delta(\{1,2\})$ be the uninformative first-stage channel with
$B(1\mid a)=t$ for every $a$.
Both branches are reached, and their posterior is $q$.  The compounds
\[
 B\triangleright(P,R),
 \qquad
 B\triangleright(Q,R)
\]
have the posterior laws on the right-hand side of
\eqref{pt:eq:public-posterior-mixture}; they differ from the displayed public
mixtures only by record labels.  Axiom~\textup{(A8)} gives the desired
biconditional directly.  Posterior-law substitution transfers it between the
compounds and the tagged mixtures.
\end{proof}

\begin{lemma}[Affine posterior-law representation]\label{pt:lem:postsep}
For every full-support $q\in\Delta(A)$ there is an affine functional
\[
 F_q:\mathcal M_q\longrightarrow\mathbb R
\]
such that, for all finite channels $P,Q$ on $A$,
\begin{equation}
 (q,P)\succeq(q,Q)
 \quad\Longleftrightarrow\quad
 F_q(\mu_{qP})\ge F_q(\mu_{qQ}).
 \label{pt:eq:fibre-representation}
\end{equation}
If $|A|\ge2$, $F_q$ is unique up to a positive affine transformation.  If
$|A|=1$, $\mathcal M_q$ is a singleton and we take $F_q\equiv0$.
\end{lemma}

\begin{proof}
By Lemma~\ref{pt:lem:plsuff} and
\eqref{pt:eq:canonical-implementation}, $\succeq_q$ is a weak order on the
mixture set $\mathcal M_q$.  Lemma~\ref{pt:lem:publiccoin} gives mixture
independence.  It remains only to verify Herstein--Milnor's segment-continuity
condition.  This is where primitive
Axiom~\textup{(A2)} is used; no topology on the collection of all finite
posterior laws is needed.

Fix $\mu,\nu,\rho\in\mathcal M_q$ and choose finite implementing channels
$P,Q,R$.  For $t\in[0,1]$ define
\[
 M_t:=tP\oplus(1-t)Q.
\]
All $M_t$ have the single fixed record alphabet given by the tagged union of
the record alphabets of $P$ and $Q$.  If $t_n\to t$, then $M_{t_n}\to M_t$
entrywise.  After adjoining the fixed block $R$, all comparisons occur in one
fixed finite action and record alphabet.  Axiom~\textup{(A2)} therefore makes
both sets
\[
 \{t\in[0,1]:M_t\succeq_q R\},
 \qquad
 \{t\in[0,1]:R\succeq_q M_t\}
\]
closed.  By \eqref{pt:eq:public-posterior-mixture}, these are exactly
\[
 \{t:t\mu+(1-t)\nu\succeq_q\rho\},
 \qquad
 \{t:\rho\succeq_qt\mu+(1-t)\nu\}.
\]
Thus the Herstein--Milnor mixture-space theorem
\cite{HersteinMilnor1953} applies and supplies an affine representation
$F_q$.  When $|A|\ge2$, Axiom~\textup{(A4)} makes the order nonconstant, so
the usual positive-affine uniqueness conclusion applies.  For a singleton
action set, Bayes plausibility forces every posterior law to equal
$\delta_q$.
\end{proof}

The preceding proof uses only closedness along each particular public-mixture
segment.  It neither assumes nor proves closedness under arbitrary weakly
convergent sequences whose support sizes grow.  It also requires no
pointwise integrand on $\Delta(A)$: every subsequent calculation uses only the
affinity of $F_q$ on $\mathcal M_q$ and its linear part on the affine hull.

Let
\[
 \delta_q\in\mathcal M_q
\]
be the posterior law of the uninformative channel and let
\[
 \chi_q:=\sum_{a\in A}q(a)\delta_{e_a}\in\mathcal M_q
\]
be the posterior law of full revelation.

\begin{lemma}[Canonical fibre normalization]\label{pt:lem:convnorm}
The representatives in Lemma~\ref{pt:lem:postsep} may be chosen so that
\begin{equation}
 F_q(\delta_q)=0,
 \qquad
 F_q(\mu)\ge0\quad(\mu\in\mathcal M_q),
 \label{pt:eq:zero-minimum}
\end{equation}
and, whenever $|A|\ge2$,
\begin{equation}
 F_q(\chi_q)=1.
 \label{pt:eq:full-revelation-anchor}
\end{equation}
These two anchors select a unique affine representative on every
nondegenerate full-support fibre.  On a singleton fibre we retain
$F_q\equiv0$.
\end{lemma}

\begin{proof}
Every channel $P$ can be record-postprocessed to the uninformative channel by
deleting its record.  Axiom~\textup{(A6)} and
\eqref{pt:eq:fibre-representation} therefore imply
\[
 F_q(\mu_{qP})\ge F_q(\delta_q).
\]
Every law in $\mathcal M_q$ is implemented by
\eqref{pt:eq:canonical-implementation}, so $\delta_q$ is a minimum of $F_q$
on the whole fibre.  Subtract $F_q(\delta_q)$ to obtain
\eqref{pt:eq:zero-minimum}.

If $|A|\ge2$, Axiom~\textup{(A4)} says that full revelation is strictly
preferred to no information at every full-support prior.  Hence
$F_q(\chi_q)>0$.  Dividing by this positive number gives
\eqref{pt:eq:full-revelation-anchor} without changing any comparison.  If two
affine representatives have both anchor values, positive-affine uniqueness
forces the additive constant to be zero and the positive multiplier to be
one.  The singleton convention is immediate.
\end{proof}

\subsubsection{Stage 2: transport across priors and support faces}

The remaining front-end issue is transport.  The same argument handles
renaming actions, deleting distinctions that the channel cannot detect, and
passing to the positive support of a boundary prior.  Keeping these facts in
one lemma makes clear that all three are consequences of action processing,
not extra invariance conventions.

For a bijection $\pi:A\to A'$, write $q\pi$ for the transported prior and
$\pi_\#\mu$ for the posterior law obtained by transporting every atom of
$\mu$.  If $B=\operatorname{supp}(q)$, write $q_B$ for the full-support
restriction of $q$ to $B$, $P_B$ for the restricted channel, and
$\pi_B\mu$ for the law obtained by restricting every posterior atom to $B$.

\begin{lemma}[Relabelling, undetectable distinctions, and support transport]
\label{pt:lem:transport}\label{pt:lem:actionbase}\label{pt:lem:neutrality}\label{pt:lem:supprestrict}
The canonically normalized fibre values have the following properties.

\begin{enumerate}[label=\textup{(\alph*)},leftmargin=*]
\item \emph{Exact relabelling.}
If $q$ has full support and $\pi:A\to A'$ is a bijection, then
\begin{equation}
 F_{q\pi}^{A'}(\pi_\#\mu)=F_q^A(\mu)
 \qquad(\mu\in\mathcal M_q).
 \label{pt:eq:exact-relabel}
\end{equation}

\item \emph{Undetectable distinctions.}
Let $S:A\to B$ be an onto deterministic map, let $q$ have full support, and
suppose that $P:A\to\Delta(X)$ is constant on every fibre of $S$.  Then
\begin{equation}
 (q,P)\sim(qS,S^qP).
 \label{pt:eq:undetectable}
\end{equation}
In particular, if $G_S$ reports only $S(a)$, then
$(q,G_S)\sim(qS,\Id_B)$.

\item \emph{Support restriction.}
For an arbitrary prior $q\in\Delta(A)$ with nonempty support
$B:=\operatorname{supp}(q)$ and every channel $P$,
\begin{equation}
 (q,P)\sim(q_B,P_B).
 \label{pt:eq:support-indifference}
\end{equation}
Consequently comparisons at $q$ depend only on the rows indexed by $B$, and
\begin{equation}
 (q,P)\succeq(q,Q)
 \quad\Longleftrightarrow\quad
 (q_B,P_B)\succeq(q_B,Q_B).
 \label{pt:eq:support-comparison}
\end{equation}
Every atom of a law $\mu\in\mathcal M_q$ is supported on $B$, and the
definition
\begin{equation}
 F_q^A(\mu):=F_{q_B}^{B}(\pi_B\mu)
 \label{pt:eq:boundary-representative}
\end{equation}
represents all continuation comparisons at the boundary prior $q$.
\end{enumerate}
\end{lemma}

\begin{proof}
\emph{Part (a).}
Regard $\pi$ as a deterministic action processor.  Axiom~\textup{(A7)}
gives
\[
 (q,P)\succeq(q\pi,P\pi^{-1}),
\]
where $(P\pi^{-1})(x\mid a')=P(x\mid\pi^{-1}(a'))$.  Applying the same axiom
to $\pi^{-1}$ gives the reverse comparison, so the two pairs are indifferent.
Apply this simultaneously to channels implementing arbitrary
$\mu,\nu\in\mathcal M_q$ and use Lemma~\ref{pt:lem:blockcoh} in a common
four-block environment.  One obtains
\[
 F_q^A(\mu)\ge F_q^A(\nu)
 \quad\Longleftrightarrow\quad
 F_{q\pi}^{A'}(\pi_\#\mu)
 \ge F_{q\pi}^{A'}(\pi_\#\nu).
\]
Thus $G(\mu):=F_{q\pi}^{A'}(\pi_\#\mu)$ is an affine representation of the
same nonconstant order as $F_q^A$.  Relabelling sends
$\delta_q$ to $\delta_{q\pi}$ and $\chi_q$ to $\chi_{q\pi}$.  The two anchor
equalities in Lemma~\ref{pt:lem:convnorm}, followed by affine uniqueness, give
$G=F_q^A$.  If $A$ is a singleton, both sides vanish.

\emph{Part (b).}
Since $q$ has full support and $S$ is onto, $qS$ has full support.  One
direction of \eqref{pt:eq:undetectable} is Axiom~\textup{(A7)}.  For the
reverse direction, define the Bayesian refinement kernel
\[
 R(a\mid b):=
 \begin{cases}
 q(a)/(qS)(b),&S(a)=b,\\
 0,&S(a)\ne b.
 \end{cases}
\]
Then $(qS)R=q$.  Since $P$ is constant on each fibre of $S$, the pushforward
$S^qP$ has at $b$ precisely the common row of $P$ on $S^{-1}(b)$, and pushing
it back through $R$ recovers $P$.  Axiom~\textup{(A7)} applied to
$(qS,S^qP)$ and $R$ therefore gives the reverse weak comparison.  For the
last assertion, $G_S$ is $S$-measurable and its pushforward is $\Id_B$.

\emph{Part (c).}
Choose $b_0\in B$ and let $S:A\to B$ fix every $b\in B$ and send every
$a\notin B$ to $b_0$.  Since $q$ assigns zero mass outside $B$,
\[
 qS=q_B,
 \qquad
 S^qP=P_B.
\]
Axiom~\textup{(A7)} gives $(q,P)\succeq(q_B,P_B)$.  Conversely, let
$\iota:B\hookrightarrow A$ be the inclusion.  Then $q_B\iota=q$.  The exact
Bayesian completion equation pins down the processed channel on $B$ and leaves
all zero-mass rows in $A\setminus B$ arbitrary.  The all-completions clause of
Axiom~\textup{(A7)} therefore permits us to choose the completion to be the
original $P$, yielding $(q_B,P_B)\succeq(q,P)$.  This proves
\eqref{pt:eq:support-indifference}; applying it to $P$ and $Q$ in a common
block environment proves \eqref{pt:eq:support-comparison}.  It also shows
immediately that channels agreeing on $B$ are indifferent at $q$.

Finally, if $\mu=\sum_k\lambda_k\delta_{s_k}\in\mathcal M_q$ with
$\lambda_k>0$, then for $a\notin B$,
\[
 0=q(a)=\sum_k\lambda_ks_k(a).
\]
Nonnegativity forces $s_k(a)=0$ for every $k$.  Thus every atom lies on the
face $\Delta(B)$ and $\pi_B\mu\in\mathcal M_{q_B}$.  The comparison
equivalence just proved and the full-support representation on $B$ show that
\eqref{pt:eq:boundary-representative} represents precisely the boundary
continuation order.  This definition is also compatible with relabelling by
part~(a).
\end{proof}

At this point every later calculation may work entirely with the affine
functionals $F_q$, their linear parts on the affine hulls of
$\mathcal M_q$, and the exact transport identities above.  No global
posterior-law continuity statement and no posterior-separable integrand are
needed.  From now on ambient action-set superscripts are suppressed unless
two action sets occur in the same formula; a boundary-prior subscript always
means the representative transported from its positive support.  The next two
stages derive product and branch scale coherence, with Axiom~\textup{(A8)} and its finite-branch extension providing
the essential cardinal alignment across reached
branches.

\subsubsection{Stage 3: products and a common comparison scale}

The value $F_q$ compares experiments only within the fibre of one prior.
Embedding experiments from two priors into their product fibre will put both
values on one cardinal comparison scale.  We first derive the only
separability fact about independent products needed for that construction.
Write $P\succeq_qQ$ for $(q,P)\succeq(q,Q)$, and similarly for strict
preference.  For a channel $P:A\to\Delta(X)$, let $\widetilde P$ denote its
lift to $A\times B$, constant in the $B$-coordinate; for
$R:B\to\Delta(Y)$, let $\widetilde R$ denote the symmetric lift, constant in
the $A$-coordinate.
Posterior laws may be substituted throughout by Lemma~\ref{pt:lem:plsuff};
in particular, adding null records and permuting record labels has no effect
on a comparison.

\begin{lemma}[Cancellation of a common independent background]
\label{pt:lem:background}
Let $p\in\Delta(A)$ and $r\in\Delta(B)$ have full support.  For channels
$P,Q$ on $A$ and a channel $R$ on $B$,
\begin{equation}
 P\otimes R\succeq_{p\otimes r}Q\otimes R
 \quad\Longleftrightarrow\quad
 P\succeq_p Q .
\tag{BI}\label{pt:eq:backgroundinert}
\end{equation}
The symmetric assertion holds with the two coordinates interchanged.  The
equivalence, including its strict form, remains valid when a posterior reached
in the background coordinate lies on a proper support face.
\end{lemma}

\begin{proof}
Pad $P$ and $Q$ by null records so that they have the same record alphabet.
Regard the lift $\widetilde R$ of $R$ as a first-stage experiment on
$A\times B$, followed on every reached branch by $\widetilde P$, respectively
$\widetilde Q$.  Up to the harmless permutation of the two record
coordinates, the resulting compounds are $P\otimes R$ and $Q\otimes R$.

If $y$ is reached under $R$, its posterior on $A\times B$ is
\[
 p\otimes r_y,
 \qquad r_y:=r_y^{\,r,R}.
\]
The coordinate projection $A\times\operatorname{supp}(r_y)\to A$ and the
support-transport clauses of Lemmas~\ref{pt:lem:neutrality} and
\ref{pt:lem:supprestrict} give
\[
 (p\otimes r_y,\widetilde P)\sim(p,P),
 \qquad
 (p\otimes r_y,\widetilde Q)\sim(p,Q).
\tag{BI$'$}\label{pt:eq:bg-branch-transport}
\]
Thus every reached branch has the same weak comparison as $P$ versus $Q$ at
$p$.  If $P\succeq_pQ$, Lemma~\ref{pt:lem:finite-branch-extension} applied to
these branches proves the forward implication in~\eqref{pt:eq:backgroundinert}.

For the converse, suppose that $P\not\succeq_pQ$.  Completeness gives
$Q\succ_pP$.  By~\eqref{pt:eq:bg-branch-transport}, the reversed comparison is
strict on every reached branch; at least one such branch has positive mass.
The strict part of Lemma~\ref{pt:lem:finite-branch-extension} therefore gives
$Q\otimes R\succ_{p\otimes r}P\otimes R$, contradicting
$P\otimes R\succeq_{p\otimes r}Q\otimes R$.  This proves the converse and
also records why weak branch monotonicity alone would not suffice.  Interchange
the coordinates for the symmetric assertion.
\end{proof}

For finite posterior laws define
\[
 \left(\sum_i m_i\delta_{u_i}\right)\otimes
 \left(\sum_j n_j\delta_{v_j}\right)
 :=\sum_{i,j}m_i n_j\delta_{u_i\otimes v_j}.
\]
Bayes' rule gives
\begin{equation}
 \mu_{p\otimes r,P\otimes R}=\mu_{pP}\otimes\mu_{rR}.
\tag{PL}\label{pt:eq:productlaw}
\end{equation}

\begin{lemma}[Product gauge]
\label{pt:lem:coherentnorm}
There are positive rescalings $G_q=c_qF_q$ of the zero-normalised affine
representatives and a single number $\kappa\in\mathbb R$ such that
\begin{equation}
 G_{p\otimes r}(\mu\otimes\nu)
 =G_p(\mu)+G_r(\nu)+\kappa G_p(\mu)G_r(\nu)
\tag{QA}\label{pt:eq:QA}
\end{equation}
for all full-support priors $p,r$ and all
$\mu\in\mathcal M_p$, $\nu\in\mathcal M_r$.  These rescalings can be chosen
exactly invariant under bijections.  If a factor is a singleton, its
representative is zero and the same formula holds.  Finally,
\begin{equation}
 1+\kappa G_r(\nu)>0
\tag{POS}\label{pt:eq:QAslope}
\end{equation}
whenever the $r$-fibre is nonconstant.
\end{lemma}

\begin{proof}
We give the algebra because both the common value of $\kappa$ and the strict
inequality~\eqref{pt:eq:QAslope} are used later.  Initially retain the notation
$F_q$.  For nondegenerate $p,r$, put
\[
 L_{p,r}(\mu,\nu):=F_{p\otimes r}(\mu\otimes\nu),
 \qquad x:=F_p(\mu),\quad y:=F_r(\nu).
\]
The map $L_{p,r}$ is separately affine.  By
Lemma~\ref{pt:lem:background}, every first-coordinate slice represents the
$F_p$-order and every second-coordinate slice represents the $F_r$-order.
Affine-utility uniqueness, first on a slice and then as the other coordinate
varies affinely, therefore gives unique constants
$A_{p,r},B_{p,r}>0$ and $C_{p,r}\in\mathbb R$ such that
\begin{equation}
 L_{p,r}(\mu,\nu)
 =A_{p,r}x+B_{p,r}y+C_{p,r}xy.
\label{pt:eq:pairbilinear}
\end{equation}
Here is the coefficient argument.  Fixing $\nu$ first gives
$L_{p,r}=\alpha(\nu)x+\gamma(\nu)$ with $\alpha(\nu)>0$.  At
$\mu=\delta_p$, the function
$\gamma(\nu)=L_{p,r}(\delta_p,\nu)$ is an affine representative of the
$F_r$-order and vanishes at $\delta_r$; hence
$\gamma(\nu)=B_{p,r}y$ for some $B_{p,r}>0$.  Choose
$\bar\mu$ with $\bar x:=F_p(\bar\mu)\ne0$.  The other slice has the form
\[
 L_{p,r}(\bar\mu,\nu)=A_{p,r}\bar x+D_{p,r}y,
\]
where $A_{p,r}>0$ is its value slope at $\nu=\delta_r$.  Comparing this with
$\alpha(\nu)\bar x+B_{p,r}y$ shows that
$\alpha(\nu)=A_{p,r}+C_{p,r}y$, where
$C_{p,r}:=(D_{p,r}-B_{p,r})/\bar x$.  This proves
\eqref{pt:eq:pairbilinear}; the symmetric calculation shows that the two
slice slopes obey
\begin{equation}
 A_{p,r}+C_{p,r}y>0,
 \qquad B_{p,r}+C_{p,r}x>0
\label{pt:eq:pair-slopes}
\end{equation}
throughout the two value ranges.

Exact relabelling makes the evaluations under the two associations of a
triple product equal.  Comparing the coefficients of $x,y,z$ gives
\begin{align}
 A_{p\otimes r,s}A_{p,r}&=A_{p,r\otimes s},\label{pt:eq:Acoc1}\\
 A_{p\otimes r,s}B_{p,r}&=B_{p,r\otimes s}A_{r,s},\label{pt:eq:Acoc2}\\
 B_{p\otimes r,s}&=B_{p,r\otimes s}B_{r,s}.
\label{pt:eq:Acoc3}
\end{align}
The coordinate swap similarly gives
\[
 B_{r,p}=A_{p,r},\qquad A_{r,p}=B_{p,r},\qquad C_{r,p}=C_{p,r}.
\tag{SW}\label{pt:eq:coefficient-swap}
\]
This coefficient comparison is legitimate: the laws
$(1-t)\delta_p+t\chi_p$ make $x$ range over a nondegenerate interval, and
the same holds for $y$ and $z$.
Let $\rho(p,r):=B_{p,r}/A_{p,r}$.  Dividing
\eqref{pt:eq:Acoc2} by~\eqref{pt:eq:Acoc1} yields
\[
 \rho(p,r)=\rho(p,r\otimes s)A_{r,s}.
\]
Fix a nondegenerate reference prior $q_0$ and set
$g(r):=\rho(q_0,r)$.  Then
\begin{equation}
 A_{r,s}=\frac{g(r)}{g(r\otimes s)}.
\label{pt:eq:gaugeA}
\end{equation}
The uniqueness of the coefficients in~\eqref{pt:eq:pairbilinear}, applied
after product relabellings, shows that $g$ is invariant under bijections.
Consequently the positive rescaling $G_p:=g(p)F_p$ preserves exact
relabelling.  Equation~\eqref{pt:eq:gaugeA} makes the new first coefficient
equal to one, and~\eqref{pt:eq:coefficient-swap} makes the second coefficient
equal to one as well.  We henceforth use $G$ for these gauged representatives.

Write the remaining coefficient as $\kappa_{p,r}$.  Associativity of
\[
 x\oplus_{p,r}y:=x+y+\kappa_{p,r}xy
\]
and comparison of the coefficients of $xy,xz,yz$ give
\[
 \kappa_{p,r}=\kappa_{p,r\otimes s},\qquad
 \kappa_{p\otimes r,s}=\kappa_{p,r\otimes s},\qquad
 \kappa_{p\otimes r,s}=\kappa_{r,s}.
\tag{KA}\label{pt:eq:kappa-assoc}
\]
(The $xyz$ coefficient gives the product of the corresponding two sides and
adds no restriction.)  Taking $s=q_0$ in~\eqref{pt:eq:kappa-assoc} shows that
$\kappa_{p,r}=\kappa_{r,q_0}$, so it is independent of $p$; the swap in
\eqref{pt:eq:coefficient-swap} makes it symmetric, hence independent of $r$.
Call the common value $\kappa$.  After the gauge,
\eqref{pt:eq:pair-slopes} is exactly
$1+\kappa G_r(\nu)>0$ and its symmetric counterpart.

For a singleton prior set $G\equiv0$; the canonical bijections
$A\cong A\times\{*\}$ give~\eqref{pt:eq:QA}.  For a boundary prior, define
$G_r$ by exact transport from its positive-probability support, as in
Lemma~\ref{pt:lem:supprestrict}.  This convention is relabelling invariant.
\end{proof}

\begin{corollary}[Exact relabelling after the product gauge]
\label{pt:cor:permutationinvariance}
For every bijection $\pi:A\to A'$, prior $q\in\Delta(A)$, and
$\mu\in\mathcal M_q$,
\[
 G_{q\pi}^{A'}(\pi_\#\mu)=G_q^A(\mu).
\]
\end{corollary}

\begin{proof}
The initial equality is Lemma~\ref{pt:lem:actionbase}; the proof of
Lemma~\ref{pt:lem:coherentnorm} shows that its gauge factor is invariant under
bijections.
\end{proof}

\begin{lemma}[Cross-prior block representation]
\label{pt:lem:blockbridge}
For arbitrary finite priors $q\in\Delta(A)$ and $r\in\Delta(B)$ and channels
$P,Q$ on the respective action sets,
\begin{equation}
 q^0\succeq_{P\sqcup Q}r^1
 \quad\Longleftrightarrow\quad
 G_q(\mu_{qP})\ge G_r(\mu_{r,Q}).
\tag{BC}\label{pt:eq:blockcomp}
\end{equation}
\end{lemma}

\begin{proof}
First suppose that $q,r$ have full support.  By~\eqref{pt:eq:QA} and
$G_q(\delta_q)=G_r(\delta_r)=0$,
\[
 G_{q\otimes r}(\mu_{qP}\otimes\delta_r)=G_q(\mu_{qP}),
 \qquad
 G_{q\otimes r}(\delta_q\otimes\mu_{r,Q})=G_r(\mu_{r,Q}).
\]
The two laws on the left lie in the same fibre and are implemented by
$P\otimes U_B$ and $U_A\otimes Q$, so the affine fibre representation
compares them by the displayed numbers.

It remains only to remove the uninformative product coordinates.  Projection
$A\times B\to A$ and Axiom~\textup{(A7)} give
$(q\otimes r,P\otimes U_B)\succeq(q,P)$.  Conversely the stochastic
embedding
\[
 T((a,b)\mid a')=\mathbf 1_{\{a=a'\}}r(b)
\]
sends $q$ to $q\otimes r$ and its Bayesian pushforward sends $P$ to
$P\otimes U_B$; Axiom~\textup{(A7)} gives the reverse comparison.  Thus the
two pairs are indifferent.  The same argument identifies
$(q\otimes r,U_A\otimes Q)$ with $(r,Q)$.  Block coherence and transitivity
now prove~\eqref{pt:eq:blockcomp}.  Arbitrary priors follow by restricting
both pairs and both representatives to their supports; singleton supports use
the zero convention.  The singleton record tags created by these projections
and embeddings are the canonical relabellings covered by the convention at
the start of the appendix.
\end{proof}

\subsubsection{Stage 4: sequential decomposition and coherent face scales}

Changing a continuation on one reached branch changes the posterior law by a
signed measure with zero total mass and zero barycentre.  The following space
is exactly the space of such feasible directions.
For a nonempty finite set $A$, let
\[
 W_A:=\left\{\eta:
 \begin{array}{l}
 \eta\text{ is a finitely supported signed measure on }\Delta(A),\\[-2pt]
 \eta(\Delta(A))=0,\quad \displaystyle\int p\,d\eta(p)=0
 \end{array}\right\}.
\]
If $\iota:B\hookrightarrow A$ is an injection, let
$i_\iota:W_B\to W_A$ push every posterior forward along $\iota$; for a
literal subset $B\subseteq A$, write $i_B$ for the inclusion case.  Let
$L_q^A$ be the linear part of the affine functional $G_q$ on the affine hull
of $\mathcal M_q$.

\begin{lemma}[Branch aggregation]
\label{pt:lem:branchagg}
Let $q$ have full support on $A$.  For every reached posterior $r$ with
$B:=\operatorname{supp}(r)$ and $|B|\ge2$, there is a unique
$\beta_A(q,r)>0$ satisfying
\begin{equation}
 L_q^A\circ i_B=\beta_A(q,r)L_{r_B}^B
 \quad\hbox{on }W_B.
\label{pt:eq:branch-linear}
\end{equation}
It depends only on $q,r$ and the face inclusion, not on the experiment that
reaches $r$.  If a first-stage channel $P_1:A\to\Delta(\{1,\ldots,k\})$ has branch masses
$m_i$ and posteriors $r_i$, then every list of continuation channels $Q_1,\ldots,Q_k$ satisfies
\begin{equation}
 G_q\bigl(\mu_{q,P_1\triangleright(Q_1,\ldots,Q_k)}\bigr)
 =G_q(\mu_{qP_1})+
 \sum_{i:m_i>0}m_i\beta_A(q,r_i)
 G_{r_i}(\mu_{r_i,Q_i}).
\tag{BA}\label{pt:eq:branchaggregation}
\end{equation}
For singleton $r_i$, the continuation value is zero and the corresponding
coefficient may be fixed by any positive convention.
\end{lemma}

\begin{proof}
We first record the common tangent space.  If $q$ has full support, then
\begin{equation}
 \operatorname{span}\{\mu-\nu:\mu,\nu\in\mathcal M_q\}=W_A.
\label{pt:eq:tangent-space}
\end{equation}
Only the reverse inclusion needs proof.  For $0\ne\eta\in W_A$, write its
Jordan decomposition as $\eta=\eta^+-\eta^-$, where both positive parts have
the same mass $c>0$ and, after division by $c$, the same barycentre $b$.
Choose $\varepsilon>0$ small enough that
$\tau=(q-\varepsilon b)/(1-\varepsilon)$ belongs to $\Delta(A)$.  Then
\[
 \mu^\pm:=\frac{\varepsilon}{c}\eta^\pm
          +(1-\varepsilon)\delta_\tau
 \in\mathcal M_q,
 \qquad
 \eta=\frac c\varepsilon(\mu^+-\mu^-),
\]
which proves~\eqref{pt:eq:tangent-space}.  The same argument intrinsically on
a face gives the space $W_B$.

Fix a nondegenerate $r$ supported on $B$.  It is reachable from $q$: choose
\[
 0<\varepsilon<\min\left\{1,
       \min_{b\in B}\frac{q(b)}{r(b)}\right\}
\]
and use a two-record experiment whose first record has conditional probability
$\varepsilon r(a)/q(a)$.  That record has mass $\varepsilon$ and posterior
$r$.

Implement $\nu,\nu'\in\mathcal M_{r_B}$ by channels on $B$, extend their
off-support rows arbitrarily to $A$, and pad their record alphabets if
necessary.  Vary only the continuation on the branch reaching $r$ from
$\nu$ to $\nu'$.  The aggregate posterior law changes by
$\varepsilon i_B(\nu-\nu')$.  Corollary~\ref{pt:cor:one-branch-a8}, support transport, and strictness give all
three signs:
\begin{align*}
 L_{r_B}^B(\nu-\nu')>0&\Longrightarrow
     L_q^A i_B(\nu-\nu')>0,\\
 L_{r_B}^B(\nu-\nu')<0&\Longrightarrow
     L_q^A i_B(\nu-\nu')<0,\\
 L_{r_B}^B(\nu-\nu')=0&\Longrightarrow
     L_q^A i_B(\nu-\nu')=0.
\end{align*}
The second line is the first applied after swapping the continuations.  For
the third, indifference supplies both weak branch comparisons, and
Corollary~\ref{pt:cor:one-branch-a8} supplies both aggregate inequalities.  By
\eqref{pt:eq:tangent-space}, the sign agreement extends by positive scaling
to all of $W_B$.  Positive trace orientation makes
\[
 L_{r_B}^B(\chi_{r_B}-\delta_{r_B})
 =G_{r_B}(\chi_{r_B})>0,
\]
so $L_{r_B}^B\ne0$.  Two nonzero linear functionals with the same positive
half-space have the same kernel and are positive scalar multiples; this gives
the unique coefficient in~\eqref{pt:eq:branch-linear}.  Since both linear
functionals in that equation are fixed before a reaching experiment is
chosen, the coefficient is path-independent.

For the aggregation formula, put $B_i:=\operatorname{supp}(r_i)$ and
\[
 \nu_i:=\pi_{B_i}\mu_{r_i,Q_i}\in\mathcal M_{(r_i)_{B_i}}.
\]
The compound law differs from the first-stage law by
\[
 \mu_{q,P_1\triangleright(Q_1,\ldots,Q_k)}-\mu_{qP_1}
 =\sum_{i:m_i>0}m_i\,i_{B_i}
   (\nu_i-\delta_{(r_i)_{B_i}}).
\]
Apply $L_q^A$, use~\eqref{pt:eq:branch-linear}, and use
$G_{r_i}(\delta_{r_i})=0$.  Singleton summands vanish.  This is
\eqref{pt:eq:branchaggregation}.
\end{proof}

For an injection $\iota:B\hookrightarrow A$ and full-support
$r\in\Delta(B)$, define
\[
 \beta_\iota(q,r):=\beta_A(q,\iota_\#r).
\]
The image posterior $\iota_\#r$ is reachable by the two-record construction
in the proof, and its support inclusion is precisely $i_\iota$.  Thus
\[
 L_q^A i_\iota=\beta_\iota(q,r)L_r^B
 \quad\hbox{on }W_B.
\]

\begin{lemma}[Cocycle, face coherence, and the chain gauge]
\label{pt:lem:chain}
The coefficients of Lemma~\ref{pt:lem:branchagg} admit positive scales
$a_q^A$ which are invariant under relabelling and satisfy
\begin{equation}
 \beta_\iota(q,r)=\frac{a_q^A}{a_r^B}
\tag{FS}\label{pt:eq:facescale}
\end{equation}
whenever $\iota:B\hookrightarrow A$, both sets are nondegenerate, $q$ and
$r$ have full support on their respective sets, and $\beta_\iota(q,r)$ is the
coefficient for reaching the face $\iota(B)$.  Hence the representatives
\[
 V_q:=\frac{G_q}{a_q}
\]
satisfy, for every $P_1:A\to\Delta(\{1,\ldots,k\})$ with branch masses $m_i$ and posteriors $r_i$, and every
list of continuation channels $Q_1,\ldots,Q_k$, the exact chain rule
\begin{equation}
 V_q\bigl(\mu_{q,P_1\triangleright(Q_1,\ldots,Q_k)}\bigr)
 =V_q(\mu_{qP_1})
  +\sum_{i:m_i>0}m_iV_{r_i}(\mu_{r_i,Q_i}).
\tag{CR}\label{pt:eq:chainrule}
\end{equation}
Singleton continuation terms are zero.
\end{lemma}

\begin{proof}
There are two normalization tasks: first solve the branch cocycle within each
fixed action-set dimension, then remove the residual scalar attached to
embedding an $m$-point face in an $n$-point simplex.

We first prove the cocycle before choosing scales.  For nested injections
$C\xhookrightarrow{\jmath}B\xhookrightarrow{\iota}A$ and full-support
priors $q,r,s$ on $A,B,C$, respectively, the defining identities on $W_C$
are
\[
 L_q^A i_{\iota\circ\jmath}
 =\beta_\iota(q,r)L_r^B i_\jmath
 =\beta_\iota(q,r)\beta_\jmath(r,s)L_s^C.
\]
The direct coefficient from $q$ to $s$ is unique because $L_s^C\ne0$ by
Axiom~\textup{(A4)}.  Therefore
\begin{equation}
 \beta_{\iota\circ\jmath}(q,s)
 =\beta_\iota(q,r)\beta_\jmath(r,s).
\label{pt:eq:branch-cocycle}
\end{equation}
This single identity covers full-support paths, full-to-face paths, and
nested support faces.

For a fixed nondegenerate $A$, apply~\eqref{pt:eq:branch-cocycle} to identity
embeddings.  Fix a full-support basepoint $q_A$ and put
\[
 a_q^A:=\beta_{\mathrm{id}_A}(q,q_A).
\]
Then
\[
 \beta_{\mathrm{id}_A}(q,r)=\frac{a_q^A}{a_r^A}.
\tag{WS}\label{pt:eq:within-scale}
\]
Choosing $q_A$ to be uniform and using exact relabelling of $G$ makes these
initial scales invariant under bijections.  They remain free up to one
positive multiplier for each cardinality.

For an injection $\iota:B\hookrightarrow A$, define its residual defect by
\[
 K(\iota):=\beta_\iota(q,r)\frac{a_r^B}{a_q^A}.
\]
Equation~\eqref{pt:eq:within-scale} shows that this number is independent of
$q$ and $r$: changing either prior merely inserts and cancels the
corresponding within-set coefficient.  Equation~\eqref{pt:eq:branch-cocycle}
gives
\[
 K(\iota\circ\jmath)=K(\iota)K(\jmath).
\]
Exact relabelling implies that $K(\iota)$ depends only on
$n=|A|$ and $m=|B|$; write it $K_{n,m}$.  Thus
\[
 K_{n,\ell}=K_{n,m}K_{m,\ell},\qquad K_{n,n}=1.
\]
Set $t_n:=K_{n,2}$ and $t_2:=1$.  Then
$K_{n,m}=t_n/t_m$.  Replacing every $a_q^A$ on an $n$-point set by
$t_n a_q^A$ changes the defect to
$K_{n,m}t_m/t_n=1$, while preserving~\eqref{pt:eq:within-scale} and
relabel invariance.  This proves~\eqref{pt:eq:facescale}.

For a boundary prior, $a_r$ means the scale on its support; the preceding
construction says precisely that this agrees with the scale obtained by
reaching that support as a face.  Divide~\eqref{pt:eq:branchaggregation} by
$a_q$ and use~\eqref{pt:eq:facescale}.  Every nondegenerate branch coefficient
becomes
\[
 \frac1{a_q}\,\beta_A(q,r_i)G_{r_i}
 =\frac1{a_{r_i}}G_{r_i}=V_{r_i}.
\]
Singleton values vanish, so their scales and coefficients may be fixed
arbitrarily.  This proves~\eqref{pt:eq:chainrule}.
\end{proof}

\subsubsection{Stage 5: compatibility and entropy identification}

The product and sequential evaluations of full revelation now meet.  Their
compatibility first removes the provisional product interaction and then
leaves a strongly additive uncertainty functional.

\begin{lemma}[Interaction collapse and universal chain scale]
\label{pt:lem:scalecoherence}
The coefficient $\kappa$ in~\eqref{pt:eq:QA} is zero.  Moreover there is a
single $C>0$ such that $a_q=C$ for every nondegenerate finite prior, read on
its support.  The same value may be assigned to singleton priors.
Consequently the final representatives $V_q=G_q/C$ satisfy both
\begin{align}
 V_{p\otimes r}(\mu\otimes\nu)
   &=V_p(\mu)+V_r(\nu),\tag{PA}\label{pt:eq:PArecovered}\\
 V_q\bigl(\mu_{q,P_1\triangleright(Q_1,\ldots,Q_k)}\bigr)
   &=V_q(\mu_{qP_1})+\sum_{i:m_i>0}m_i
     V_{r_i}(\mu_{r_i,Q_i}).\tag{CR$'$}\label{pt:eq:final-chain}
\end{align}
for every $P_1:A\to\Delta(\{1,\ldots,k\})$ with branch masses $m_i$ and posteriors $r_i$, and every list of
continuation channels $Q_1,\ldots,Q_k$.
The cross-prior bridge of Lemma~\ref{pt:lem:blockbridge} is valid with $V$
in place of $G$.
\end{lemma}

\begin{proof}
In this proof only, put
\[
 H_G(q):=G_q(\chi_q),\qquad Z(q):=1+\kappa H_G(q).
\]
For nondegenerate $q$, Axiom~\textup{(A4)} and zero normalisation give
$H_G(q)>0$, while the strict slice-slope
\eqref{pt:eq:QAslope} gives
\begin{equation}
 Z(q)>0.
\label{pt:eq:Zpositive}
\end{equation}
For a singleton, $H_G=0$ and $Z=1$.

\emph{Product revelation versus sequential revelation.}
For nondegenerate full-support $p,r$, product quasi-additivity gives
\begin{equation}
 H_G(p\otimes r)
 =H_G(p)+H_G(r)+\kappa H_G(p)H_G(r).
\tag{QH}\label{pt:eq:QH}
\end{equation}
Compute the same full revelation by first revealing the $A$-coordinate and
then the $B$-coordinate.  The first-stage law is
$\chi_p\otimes\delta_r$ and has value $H_G(p)$ by~\eqref{pt:eq:QA}.
After action $a$ is revealed, the posterior is $e_a\otimes r$ on the face
$\{a\}\times B$; support restriction and relabelling identify the
continuation value with $H_G(r)$.  Every such branch has coefficient
$a_{p\otimes r}/a_r$ by~\eqref{pt:eq:facescale}.  Since the branch weights
sum to one, Lemma~\ref{pt:lem:branchagg} gives
\begin{equation}
 H_G(p\otimes r)=H_G(p)+\frac{a_{p\otimes r}}{a_r}H_G(r).
\tag{SA}\label{pt:eq:seqA}
\end{equation}
Comparing~\eqref{pt:eq:QH} and~\eqref{pt:eq:seqA}, and using
$H_G(r)>0$, yields
\[
 \frac{a_{p\otimes r}}{a_r}=Z(p).
\tag{R1}\label{pt:eq:ratio1}
\]
Revealing $B$ first gives symmetrically
\[
 \frac{a_{p\otimes r}}{a_p}=Z(r).
\tag{R2}\label{pt:eq:ratio2}
\]
Thus $a_rZ(p)=a_pZ(r)$ for every pair of nondegenerate priors, even on
different finite sets.  By~\eqref{pt:eq:Zpositive}, there is a universal
$C>0$ such that
\begin{equation}
 a_q=CZ(q).
\tag{AS}\label{pt:eq:ascaleZ}
\end{equation}
Assign $a=C$ on singleton supports.

If $\kappa=0$, then $Z\equiv1$ and the conclusion already follows.  It
remains to rule out $\kappa\ne0$.

\emph{The grouping weight.}
Let $q$ have full support and let
$\mathcal P=\{A_1,\ldots,A_m\}$ partition its action set into nonempty
cells.  Put
\[
 p_i=q(A_i),\qquad q^i=q(\,\cdot\mid A_i),\qquad
 p=(p_1,\ldots,p_m).
\]
If $C_{\mathcal P}$ reports only the cell, undetectable-distinctions
neutrality gives
\[
 (q,C_{\mathcal P})\sim(p,\Id_{\{1,\ldots,m\}}).
\]
The cross-prior bridge therefore identifies the coarse value as
\begin{equation}
 G_q(\mu_{q,C_{\mathcal P}})=H_G(p).
\tag{BG}\label{pt:eq:blockvalue}
\end{equation}
Append full revelation within each cell.  The final law is $\chi_q$, and the
branch formula together with~\eqref{pt:eq:ascaleZ} gives
\begin{equation}
 H_G(q)=H_G(p)+Z(q)\sum_i p_i\frac{H_G(q^i)}{Z(q^i)}.
\tag{GR}\label{pt:eq:groupingH}
\end{equation}
Here singleton cells contribute zero.

Set $w(q):=1/Z(q)>0$.  Because $\kappa\ne0$ in the present paragraph,
$\kappa H_G(x)=Z(x)-1$.  Substitution into
\eqref{pt:eq:groupingH} gives
\[
 Z(q)-1=Z(p)-1+
 Z(q)\sum_i p_i\left(1-\frac1{Z(q^i)}\right).
\]
Cancelling and inverting yields the positive grouping-weight equation
\begin{equation}
 w(q)=w(p)\sum_i p_iw(q^i).
\tag{W}\label{pt:eq:weight}
\end{equation}
All distributions in what follows are read on their positive-probability
supports.  For $u\otimes v$, partition by the first coordinate: the coarse
distribution is $u$ and every conditional distribution is $v$.  Thus
\eqref{pt:eq:weight} gives
\begin{equation}
 w(u\otimes v)=w(u)w(v).
\tag{M}\label{pt:eq:wmult}
\end{equation}

\emph{Two groupings force the weight to be constant.}
Take arbitrary finite probability vectors $u,v$ and write
$x=w(u)$ and $y=w(v)$.  Let $T$ be the law of a triple
$(\ell,z_1,z_2)$: $\ell$ is uniform on $\{0,1\}$, and conditional on
$\ell=0$ the pair $(z_1,z_2)$ has law $u\otimes u$, while conditional on
$\ell=1$ it has law $v\otimes v$.  In tagged notation,
\[
 T:=\tfrac12(u\otimes u)^0\sqcup\tfrac12(v\otimes v)^1.
\]
Grouping first by $\ell$ and using~\eqref{pt:eq:wmult} gives
\begin{equation}
 w(T)=w(\tfrac12,\tfrac12)\frac{x^2+y^2}{2}.
\tag{E1}\label{pt:eq:twoeval1}
\end{equation}
Instead group first by $(\ell,z_1)$.  That marginal is
$S=\tfrac12u^0\sqcup\tfrac12v^1$, so
\[
 w(S)=w(\tfrac12,\tfrac12)\frac{x+y}{2}.
\]
The conditional law of $z_2$ is $u$ when $\ell=0$ and $v$ when $\ell=1$.
A second use of~\eqref{pt:eq:weight} therefore gives
\begin{equation}
 w(T)=w(S)\frac{x+y}{2}
 =w(\tfrac12,\tfrac12)\left(\frac{x+y}{2}\right)^2.
\tag{E2}\label{pt:eq:twoeval2}
\end{equation}
The common leading factor is positive.  Comparing
\eqref{pt:eq:twoeval1} and~\eqref{pt:eq:twoeval2} yields $(x-y)^2=0$.
Thus $w$ is constant on all finite probability vectors.  Its singleton value
is one, so $w\equiv1$.

It follows that $Z\equiv1$, hence $\kappa H_G(q)=0$ for every $q$.
Axiom~\textup{(A4)} supplies a nondegenerate $q$ with $H_G(q)>0$, so
$\kappa=0$, contradicting the temporary alternative.  Equation
\eqref{pt:eq:ascaleZ} now gives $a_q=C$ universally.
Dividing the product and branch formulas by $C$ proves
\eqref{pt:eq:PArecovered} and~\eqref{pt:eq:final-chain}.  Since the division
is by one common positive number, it also preserves the cross-prior bridge.
\end{proof}

\paragraph{Entropy identification.}

\begin{lemma}[Entropy reduction and Faddeev identification]
\label{pt:lem:faddeevsketch}
Define the value of full revelation by
\[
 \mathcal H(q):=V_q(\chi_q),
\]
so $\mathcal H(q)=H_G(q)/C$, with boundary priors read on their support.
Then every finite channel
$P:A\to\Delta(X)$ satisfies
\begin{equation}
 V_q(\mu_{qP})
 =\mathcal H(q)-\sum_{x:m_{qP}(x)>0}
    m_{qP}(x)\mathcal H(r_x^{\,q,P}).
\tag{ER}\label{pt:eq:entropy-reduction}
\end{equation}
Moreover,
\[
 \mathcal H(q)=\alpha\Sh(q)
\]
for a single $\alpha>0$.  Consequently
\[
 V_q(\mu_{qP})=\alpha I_{qP}(A;X).
\tag{MI}\label{pt:eq:value-MI}
\]
\end{lemma}

\begin{proof}
\emph{Entropy reduction.}
Assume first that $q$ has full support.  Append full revelation after $P$.
The final posterior law is $\chi_q$, while the continuation after record $x$
has full-revelation law $\chi_{r_x}$ on the support of $r_x$.  The exact
chain rule~\eqref{pt:eq:final-chain} therefore gives
\[
 \mathcal H(q)=V_q(\mu_{qP})
 +\sum_{x:m_{qP}(x)>0}m_{qP}(x)\mathcal H(r_x),
\]
which is~\eqref{pt:eq:entropy-reduction}.  A boundary prior is reduced to its
support; null action rows and null records affect neither side.

\emph{Strong additivity.}
Let $\mathcal P=\{A_1,\ldots,A_m\}$ be a partition of the support of $q$,
and put
\[
 p_i=q(A_i),\qquad q^i=q(\,\cdot\mid A_i),\qquad
 p=(p_1,\ldots,p_m).
\]
The channel $C_{\mathcal P}$ which reports the cell is neutral to the
corresponding full revelation on the cell index:
\[
 (q,C_{\mathcal P})\sim(p,\Id_{\{1,\ldots,m\}})
\]
by Lemma~\ref{pt:lem:neutrality}.  The common-scale cross-prior bridge from
Lemma~\ref{pt:lem:scalecoherence} therefore gives
\[
 V_q(\mu_{q,C_{\mathcal P}})=\mathcal H(p).
\]
Applying~\eqref{pt:eq:entropy-reduction} to $C_{\mathcal P}$ yields
\begin{equation}
 \mathcal H(q)=\mathcal H(p)+\sum_i p_i\mathcal H(q^i).
\tag{FA}\label{pt:eq:faddeev-recursion}
\end{equation}
Thus $\mathcal H$ is strongly additive.  Relabelling invariance of $V$ makes
$\mathcal H$ symmetric; support restriction makes it expansive under adding
zero-probability states; and $\mathcal H(\delta_a)=0$.  Record processing
makes every experiment weakly preferable to its no-information garbling, so
$V_q(\mu)\ge0$ on every fibre and hence $\mathcal H(q)\ge0$.
Axiom~\textup{(A4)} gives strict
positivity at every nondegenerate full-support prior.
Although \eqref{pt:eq:faddeev-recursion} was derived by partitioning the
positive support, it also holds with zero outer weights: delete the null cells
by support restriction, apply the displayed recursion, and reinsert them by
expansibility.

\emph{Continuity: the binary case.}
Write $h(t):=\mathcal H(t,1-t)$.  We derive, rather than assume, its
continuity from Axiom~\textup{(A2)}.  Fix $c\ge0$.  Choose a full-support
binary prior $\theta$ with $\mathcal H(\theta)>0$, as supplied by
Axiom~\textup{(A4)}.  Choose $n\ge1$ and $\lambda\in[0,1]$ so that
\[
 c=\lambda n\mathcal H(\theta).
\]
For $\theta_n:=\theta^{\otimes n}$, product additivity
\eqref{pt:eq:PArecovered} gives
$\mathcal H(\theta_n)=n\mathcal H(\theta)$.  Let $E_{n,\lambda}$ reveal the
$n$-bit action fully with probability $\lambda$ and report a common erasure
symbol otherwise.  Its posterior law is
$\lambda\chi_{\theta_n}+(1-\lambda)\delta_{\theta_n}$, so affinity gives
\[
 V_{\theta_n}(\mu_{\theta_n,E_{n,\lambda}})=c.
\]

Now fix the block environment
\[
 K_c:=\Id_{\{0,1\}}\sqcup E_{n,\lambda}.
\]
Let $q_t^0$ put the binary prior $(t,1-t)$ on its first block and let
$\theta_n^1$ put $\theta_n$ on its second block.  The cross-prior bridge says
\begin{align*}
 h(t)\ge c&\quad\Longleftrightarrow\quad q_t^0\succeq_{K_c}\theta_n^1,\\
 h(t)\le c&\quad\Longleftrightarrow\quad \theta_n^1\succeq_{K_c}q_t^0.
\end{align*}
The environment is fixed and $t\mapsto q_t^0$ is continuous, so
Axiom~\textup{(A2)} makes both level sets closed.  When $c<0$, the superlevel
set is all of $[0,1]$ and the sublevel set is empty because $h\ge0$.
Thus every upper and lower level set of $h$ is closed, and $h$ is continuous.

\emph{Continuity on every finite simplex.}
Proceed by induction on the number of coordinates.  For
$q=(q_1,(1-q_1)\bar q)$ with $q_1<1$, recursion
\eqref{pt:eq:faddeev-recursion} for the partition
$\{\{1\},\{2,\ldots,n\}\}$ gives
\begin{equation}
 \mathcal H(q)=h(q_1)+(1-q_1)\mathcal H(\bar q).
\label{pt:eq:continuity-induction}
\end{equation}
The right side is continuous away from $q_1=1$ by binary continuity and the
induction hypothesis.  At $q_1=0$, support restriction and $h(0)=0$ give
the same formula on the face.  At $q_1=1$, the induction hypothesis makes
$\mathcal H$ bounded on the compact $(n-1)$-simplex, so the second term in
\eqref{pt:eq:continuity-induction} tends to zero and the first tends to
$h(1)=0$.  Relabelling handles the other faces.  Hence $\mathcal H$ is
continuous on each closed finite simplex.

We have now verified all hypotheses of the strong-additivity form of
Faddeev's theorem: continuity on each finite simplex, invariance under
bijections, expansibility, vanishing on point masses, and
\eqref{pt:eq:faddeev-recursion}.  The finite entropy characterisation
\cite[Theorem~6]{BaezFritzLeinster2011} therefore gives
\[
 \mathcal H(q)=\alpha\Sh(q)
\]
for some $\alpha\ge0$.  Axiom~\textup{(A4)} forces $\alpha>0$.
Substitution in~\eqref{pt:eq:entropy-reduction} gives
\[
 V_q(\mu_{qP})
 =\alpha\left(\Sh(q)-\sum_x m_{qP}(x)
                    \Sh(r_x^{\,q,P})\right)
 =\alpha I_{qP}(A;X),
\]
where the sum is over positive-mass records.  This proves
\eqref{pt:eq:value-MI}.
\end{proof}

\begin{proof}[Completion of Proposition~\ref{pt:prop:main}]
After Lemma~\ref{pt:lem:scalecoherence}, the cross-prior bridge is represented
by $V$.  Lemma~\ref{pt:lem:faddeevsketch} identifies this value as
$\alpha I$ with one $\alpha>0$.  Hence, for arbitrary pairs $(q,P)$ and
$(r,Q)$,
\[
 q^0\succeq_{P\sqcup Q}r^1
 \quad\Longleftrightarrow\quad
 I_{qP}\ge I_{r,Q},
\tag{BMI}\label{pt:eq:final-block-MI}
\]
including boundary priors by support restriction.  For one fixed channel,
the duplication clause of Axiom~\textup{(A5)} identifies
$q\succeq_Pq'$ with the left side of~\eqref{pt:eq:final-block-MI} for
$Q=P$, proving the first assertion.  For the block clause,
Axiom~\textup{(A5)} deletes all irrelevant blocks and
\eqref{pt:eq:final-block-MI} compares the two remaining pairs.  A lottery
supported on block $i$ has a constant block tag, so
\[
 I_{q_i^i,\,\bigsqcup_{k\in J}P_k}
 =I_{q_iP_i},
\]
and similarly for block $j$.  This is exactly the stated common-scale block
comparison.
\end{proof}
